\clearpage
\section{Actividades del parcial}

\begin{actividad}{Descomposición y granularidad}
Elige un proceso cotidiano diferente de los ejemplos de los apuntes.
Descompónlo primero en cinco macro-pasos y después convierte uno de ellos en
al menos seis micro-pasos precisos, finitos y sin ambigüedad.
\espaciorespuesta[6cm]
\end{actividad}

\clearpage
\begin{actividad}{Variables y tipos de datos}
Diseña las variables necesarias para registrar a un estudiante: nombre, edad,
promedio, número de materias y estado de inscripción. Indica identificador,
tipo de dato, valor de ejemplo y si debe ser variable o constante. Usa
\texttt{snake\_case} o \texttt{camelCase} de manera consistente.
\espaciorespuesta[6cm]
\end{actividad}

\clearpage
\begin{actividad}{Entrada, proceso y salida}
Escribe el pseudocódigo y dibuja el diagrama de flujo para calcular el precio
final de un producto después de aplicar un descuento porcentual. Identifica
claramente entradas, proceso y salida. Utiliza el símbolo de documento para la
salida impresa.
\espaciorespuesta[8cm]
\end{actividad}

\clearpage
\begin{actividad}{Decisiones y prueba de escritorio}
Construye un algoritmo que reciba una calificación y un porcentaje de
asistencia. Debe validar que ambos datos estén en rango y decidir si el alumno
aprueba. Define tus condiciones y realiza una prueba de escritorio con tres
casos: aprobado, reprobado y dato inválido.
\espaciorespuesta[9cm]
\end{actividad}

\clearpage
\begin{actividad}{Selección de ciclos}
Resuelve cada inciso usando el ciclo más apropiado y justifica la elección:
\begin{enumerate}[leftmargin=*]
  \item solicitar una contraseña hasta que sea correcta;
  \item imprimir las tablas del 1 al 10;
  \item leer temperaturas mientras el usuario no escriba \texttt{SALIR}.
\end{enumerate}
Incluye el pseudocódigo de cada solución y una prueba de escritorio breve.
\espaciorespuesta[11cm]
\end{actividad}

\clearpage
\begin{actividad}{Ejercicio integrador}
Diseña un cajero simplificado con consulta de saldo, retiro y salida. Valida
que el retiro sea positivo y no exceda el saldo; repite el menú hasta elegir
salir y muestra una secuencia final de tres fases. Entrega análisis de entradas
y salidas, pseudocódigo, diagrama de flujo y prueba de escritorio con un retiro
válido y uno rechazado.
\espaciorespuesta[12cm]
\end{actividad}
